\section{The complete terminal fourth-moment catalogue}\label{sec:terminals}
Let $\Sigma$ be a terminal arm multiset whose distinct residual list has at least two entries, with arbitrary original multiplicities. We use the following local estimate.

\begin{theorem}[Terminal supply]\label{thm:terminal}
For each integer $n\ge0$ there is an orientation of all the arms, with actual message $F^+_{\Sigma,n}$, determined by $\Sigma$ and $n$, such that
\begin{equation}\label{eq:terminalplus}
 \E\bigl[p(T^nF^+_{\Sigma,n})^4\bigr]\le1.
\end{equation}
Its complete reversal $F^-_{\Sigma,n}=\bar F^+_{\Sigma,n}$ satisfies
\begin{equation}\label{eq:terminalminus}
 \E\bigl[q(A^nF^-_{\Sigma,n})^4\bigr]\le1.
\end{equation}
The parameter $n$ is a prescribed integer determined by the original connector length.
\end{theorem}

The catalogue below proves the theorem, with analytic estimates in the indicated sections of Appendices~\ref{app:zero-moments}--\ref{app:odd-propagation}.

\subsection{Zero-step directions}
Every direction in this subsection includes the equal-pair multiplier $E_0$. For an even residual list, put
\[
 P=\prod_{i=1}^t h_{a_i}f_{b_i},\qquad
 a_1<b_1<\cdots<a_t<b_t.
\]
This direction $F^+_{\Sigma,0}=E_0P$ satisfies $\E[p(F^+_{\Sigma,0})^4]\le1$. The pair $(1,2)$ is treated directly in Section~\ref{loc:S13}. Every other even residual list also satisfies a weighted sixth-moment bound, which yields the fourth moment by H\"older against $p\mu$; the complete even catalogue, with its component ranges, is collected at the end of Section~\ref{loc:even-zero-catalogue}, with proofs in Appendix~\ref{q3:app:foundations} and Sections~\ref{loc:S2}--\ref{loc:S9}.

For an odd residual list $a<b_1<c_1<\cdots<b_t<c_t$, $t\ge1$, use
\begin{equation}\label{eq:zeroOdd}
 F^+_{\Sigma,0}=
 \begin{cases}
 E_0 qf_{b_1}h_{c_1}\displaystyle\prod_{i=2}^t h_{b_i}f_{c_i},&a=1,\\[3pt]
 \displaystyle\overline{E_0f_a\prod_{i=1}^t h_{b_i}f_{c_i}},&a\ge2.
 \end{cases}
\end{equation}
The first line is proved in Section~\ref{loc:S14}, for every number of pairs. The second line is the full reversal of the $q$-weighted fourth estimates for positive canonical odd messages, with minimum two in Section~\ref{loc:S15} and minimum at least three in Section~\ref{loc:S16}; there the equal-pair arms are reversed together with the residual arms, which is what converts the $q$ weight into the $p$ weight.

\subsection{Positive steps: even residual lists}
Write the residual list as $a_1<b_1<\cdots<a_t<b_t$, $t\ge1$. The following table gives the residual part of $F^+_{\Sigma,n}$ for every $n\ge1$; multiply each row by its actual $E_0$.

\begin{table}[htbp]
\centering
\caption{Positive-step terminal directions for even residual lists.}\label{tab:terminal-even}
\fontsize{9}{11}\selectfont
\setlength{\tabcolsep}{1pt}\noindent\begin{tabular}{@{}p{57mm}p{68mm}p{29mm}@{}}
\toprule
Residual condition & Residual message & Complete proof\\\midrule
$t=1$, $(a_1,b_1)=(1,2)$ & $qf_2$ & Section~\ref{loc:S21}\\
$t=1$, $a_1=1$, $b_1\ge3$ & $qf_{b_1}$ & Section~\ref{loc:S24}\\
$t\ge2$, $a_1=1$ & $\prod_{i=1}^t h_{a_i}f_{b_i}$ & Section~\ref{loc:S29}\\
$t=1$, $(a_1,b_1)=(2,3)$ & $h_2f_3$ & Section~\ref{loc:S23}\\
$t=1$, $a_1=2$, $b_1\ge4$ & $h_2f_{b_1}$ & Section~\ref{loc:S25}\\
$t\ge2$, $a_1=2$ & $\prod_{i=1}^t h_{a_i}f_{b_i}$ & Section~\ref{loc:S33}\\
$a_1=3$, any $t\ge1$ & $f_3h_{b_1}\prod_{i=2}^t h_{a_i}f_{b_i}$ & Sections~\ref{loc:S28}, \ref{loc:S30}\\
$a_1\ge4$, any $t\ge1$ & $\prod_{i=1}^t h_{a_i}f_{b_i}$ & Section~\ref{loc:S27}\\
\bottomrule
\end{tabular}

\end{table}

The row of minimum three uses $f_3h_{b_1}$, with the first pair reversed.

\subsection{Positive steps: odd residual lists}
Write
\[
 a<b_1<c_1<b_2<c_2<\cdots<b_t<c_t,\qquad t\ge1.
\]
Again the table gives only the residual part and covers all $n\ge1$. Here $\bar f_a=T^a\one$ is a path function, whereas $b_i$ in the ordered list is an arm length.

\begin{table}[htbp]
\centering
\caption{Positive-step terminal directions for odd residual lists.}\label{tab:terminal-odd}
\fontsize{9}{11}\selectfont
\setlength{\tabcolsep}{1pt}\noindent\begin{tabular}{@{}p{54mm}p{71mm}p{29mm}@{}}
\toprule
Residual condition & Residual message & Complete proof\\\midrule
$a=1,t=1,b_1=2,c_1=3$ & $qf_2h_3$ & Section~\ref{loc:S22}\\
$a=1,t=1,b_1=2,c_1\ge4$ & $qf_2h_{c_1}$ & Section~\ref{loc:S26}\\
$a=1,t=1,3\le b_1<c_1$ & $qf_{b_1}h_{c_1}$ & Section~\ref{loc:S31}\\
$a=1,t\ge2$ & $qf_{b_1}h_{c_1}\prod_{i=2}^t h_{b_i}f_{c_i}$ & Section~\ref{loc:S35}\\
$a=2,b_1=3$, any $t\ge1$ & $\bar f_2f_3h_{c_1}\prod_{i=2}^t h_{b_i}f_{c_i}$ & Section~\ref{loc:S34}\\
$a\ge2,b_1=4$ & $\bar f_af_4h_{c_1}\prod_{i=2}^t h_{b_i}f_{c_i}$ & Section~\ref{loc:S32}\\
$a\ge2,b_1\ge5$ & $\bar f_a\prod_{i=1}^t h_{b_i}f_{c_i}$ & Section~\ref{loc:S32}\\
\bottomrule
\end{tabular}

\end{table}

\paragraph{Completeness of the terminal classification.}
For an even list the minimum is $1,2,3$, or at least four; the first
two cases split into one pair and multiple pairs.  For an odd list of
minimum one, the cases split by the number of pairs and then by the second
length.  At minimum two the second length is three, four, or at
least five; at larger minimum it is at least four.  The inward
anchor $\bar f_a$ is an actual arm with all of its edges directed toward
the root.

\subsection{Restoration, reversal, and interface length}
The proofs in the tables are for the residual message and already allow arbitrary equal-pair factors. Since $0\le E_0\le1$ and $P\ge0$ give $T^n(E_0P)\le T^nP$, taking fourth powers and integrating against the nonnegative weight $p$ restores the paired arms without weakening the bound.

Applying \eqref{eq:terminalplus} at $-K$ interchanges $T$ with $A$, $p$ with $q$, and the selected arm message with its complete reversal, which proves \eqref{eq:terminalminus}.

For a connector of original length $L$, keeping the edge nearest the internal vertex leaves $n=L-1$ edges to integrate, so $L=1$ uses the zero-step catalogue and $L\ge2$ the positive-step one. In the empty-center star construction of Section~\ref{sec:emptystar} one terminal is propagated through all $L$ edges.

\subsection{The two internal side-table roles}
After equal pairing, an internal side table in a chain is of one of two kinds. Call it type R if its residual list is empty or has at least two entries. For an empty list take $V^+=E_0$, and otherwise use the zero-step catalogue. In either case
\begin{equation}\label{eq:R}
 \E[p(V^+)^4]\le1,\qquad V^-=\bar V^+,
 \qquad \E[q(V^-)^4]\le1.
\end{equation}
Call the table type S if it has exactly one residual arm, of length $a\ge1$. Put $C^+=E_0f_a$ and $C^-=\bar C^+$. The complete singleton estimate is
\begin{equation}\label{eq:S}
 \E[q^5(C^+)^4]\le1,\qquad \E[p^5(C^-)^4]\le1.
\end{equation}
For $a=1$, $q^5p^4=q(pq)^4\le q$ proves this directly. The length-two proof is Section~\ref{loc:S17}, and all lengths $a\ge3$ are proved in Section~\ref{loc:S19}.
